\subsection{Populations and Samples}

\content{
        \begin{center}
        \begin{minipage}{0.8\textwidth}
        ``There are three kinds of lies -- lies, damned lies, and statistics." \\
        \vspace{-3pt}\par\hfill (British Prime Minister Benjamin Disraeli)
        \end{minipage}
        \end{center}
}{% presentation
}{% nopresentation
}{% comments
}


\content{
        \begin{definition} The \textbf{population} of study is the group the collected
        data is intended to describe.
        \end{definition}

        It is important to specify the population because our conclusions make
        from the data could differ based on the population we wish to study. 
}{% presentation
\vspace{1in}
}{% nopresentation
\vspace{1in}
}{% comments
Sometimes known as the target population, since a poorly designed study could
fail to represent the true population.
}

\content{
        \begin{example} Suppose we wish to study the average amount of money
        spent on textbooks for first-year college students. The population we
        wish to study might be any of the following: 
        \begin{itemize}
        \item First year students at the University of Idaho.
        \item First year students at any university in Idaho.
        \item First year students at any community college in Idaho.
        \item First year students at any community college in the USA.
        \item Etc.
        \end{itemize}
        Maybe there is a class at the university of Idaho which requires a
        particularly expensive textbook, which could skew our conclusion about
        the population if the target population is first year students at 
        universities in the USA.
        \end{example}
}{% presentation
}{% nopresentation
}{% comments
}

\content{
        \begin{definition} If we are able to gather data on every single member
        of the population, any value computed using all the data is called a 
        \textbf{parameter}.
        \end{definition}
}{% presentation
}{% nopresentation
}{% comments
}

\content{
        \begin{definition} A survey of the entire population is called a
        \textbf{census}
        \end{definition}
        You are probably familiar with a few common censuses: the United States
        Census, and voting.
}{% presentation
\vspace{1in}
}{% nopresentation
}{% comments
voting in particular shows the difficulty in surveying the entire population:
some may not answer the survey.
}

\content{
        \begin{definition} A \textbf{sample} is a subset of the population,
        ideally one that is a good representation of the population.
        \end{definition}
}{% presentation
}{% nopresentation
}{% comments
}

\content{
        \begin{example}
        A researcher wanted to know how citizens of Tacoma felt about a 
        voter initiative.  To study  this, she goes to the Tacoma Mall 
        and randomly selects 500 shoppers and asks them their  opinion.  
        60\% indicate they are supportive of the initiative.  What is the 
        sample and  population?  Is the 60\% value a parameter or a statistic?
        \end{example}
}{% presentation
\vspace{1in}
}{% nopresentation
\vspace{1in}
}{% comments
Note that the effective population was mall shoppers, and there is no reason to
believe that this is a representative sample of the target population.
}

\subsection*{Categorizing Data}


\content{
        \begin{definition} \textbf{Categorical (qualitative) data} are pieces of
        information that allow us to classify the objects under investigation
        into categories.

        \textbf{Quantitative data} are responses that are numerical in nature
        that allow us to perform meaningful calculations.
        \end{definition}
}{% presentation
}{% nopresentation
}{% comments
}


\content{
        \begin{example} Suppose we put out a survey asking how many movies
        people have watched in the last year. What kind of data is this? 
        \end{example}
}{% presentation
\vspace{1in}
}{% nopresentation
\vspace{1in}
}{% comments
}

\content{
        \begin{example} Suppose we ask people what ZIP code they live in. What
        kind of data is this? 
        \end{example}
}{% presentation
\vspace{1in}
}{% nopresentation
\vspace{1in}
}{% comments
}

\subsection*{Sampling Methods}

\content{
        \begin{example} Suppose we are hired by a politician to determine how
        much support they would have if they decided to run for another term? 
        As always, the first step is identifying the population.\\
        The next step is to take a sample of the population. How do we do this?
        \end{example}
}{% presentation
\vspace{2in}
}{% nopresentation
\vspace{1.5in}
}{% comments
All people in area? All voters? All registered voters? All likely voters? This
last one is what is commonly used, but it can still fail.

In 1998, former professional wrestler Jesse Ventura was elected governor of
Minnesota.  Polls before the election consistently showed that he had little
chance of winning. One big factor of why this happened was because polls omitted
younger voters, since they were not considered likely voters. Since he had a lot
of support in this group, the polls were not accurate. 

This kind of situation is known as sampling bias, and it can be difficult to
avoid. In this example, the polls omitted a key group from their data, and so
the sample was not a good representation of the target population.
}

\content{
        \begin{definition} A sampling method is called \textbf{biased} is every member
        of the population doesn't have equal likelihood of being in the sample.
        \end{definition}
}{% presentation
}{% nopresentation
}{% comments
}
